\chapter{Construction~: un moteur d'appariement et un simulateur de bourse}\label{ch:mx:build-a-matching-engine-and-exchange-simulator}

Chaque tutoriel des livres qui suivent a besoin d'une plateforme~: quelque chose qui reçoit des ordres, les apparie, publie ce qui s'est passé dans un flux pouvant perdre des
paquets, et fait exactement la même chose à chaque exécution avec les mêmes entrées. Ce chapitre décrit celui que ce livre utilise depuis le
chapitre~1, \texttt{firm.exchsim}~: ce que fait une bourse, en messages~; le \omterm{def:mx:the-limit-order-book:engine}{moteur d'appariement}~; le flux et sa reprise~; les \omterm{def:mx:build-a-matching-engine-and-exchange-simulator:oe}{sessions de saisie des ordres}~; et
l'horloge qui le rend déterministe. Il le certifie ensuite~: la suite de conformité, l'identité octet par octet de ses trois moteurs, et leur vitesse sur cet
ordinateur portable.

\section{Ce que fait une bourse, en messages}

\begin{definition}[Simulateur de bourse]\label{def:mx:build-a-matching-engine-and-exchange-simulator:sim}
Un \emph{simulateur de bourse}\index{simulateur de bourse} (exchange simulator) est un programme qui se comporte envers ses participants comme une plateforme le ferait~: il accepte et accuse réception
des ordres par des \omterm{def:mx:build-a-matching-engine-and-exchange-simulator:oe}{sessions de saisie des ordres}, les apparie selon les règles de la plateforme, et publie les \omterm{def:mx:the-limit-order-book:event}{événements du carnet} qui en résultent dans un flux de données de marché
séquencé, avec les latences, les pertes et les pannes d'une connexion réelle, et avec des résultats qui ne dépendent que de ses entrées.
\end{definition}

\begin{definition}[Protocole de saisie des ordres, session de saisie des ordres, rapport d'exécution]\label{def:mx:build-a-matching-engine-and-exchange-simulator:oe}
Un \emph{protocole de saisie des ordres}\index{protocole de saisie des ordres} (order-entry protocol) est l'ensemble des messages qu'un participant utilise pour envoyer, modifier et annuler des ordres, et que la plateforme utilise
pour répondre. Une \emph{session de saisie des ordres}\index{session de saisie des ordres} (order-entry session) est une connexion authentifiée qui le porte, avec sa propre suite de messages, ses limites et
son état. Un \emph{rapport d'exécution}\index{rapport d'exécution} (execution report) est le message de la plateforme, sur cette session, au sujet d'un ordre~: accepté, exécuté (partiellement ou totalement),
annulé, rejeté, remplacé.
\end{definition}

\texttt{firm.exchsim} parle quatre protocoles binaires, tous générés à partir d'une seule table, \texttt{MESSAGES} dans le module de codage, avec une copie lisible par machine
dans \texttt{schema.json} à partir de laquelle le livre~13 construit ses codecs. Le \emph{flux} compte quatorze types de messages dans une disposition de type ITCH (les messages d'ajout, d'exécution, d'annulation, de suppression et de transaction du livre~1
octet par octet, plus ceux de remplacement, d'exécution à un prix, de croisement, de déséquilibre, d'événement système, d'action de négociation,
de répertoire des valeurs, d'en-tête et de pied d'instantané)~; la saisie des ordres en compte cinq (nouvel ordre, remplacement, annulation, annulation de masse, cotation de masse)~; les \omterm{def:mx:build-a-matching-engine-and-exchange-simulator:oe}{rapports d'exécution} six (accepté, exécuté,
annulé, rejeté, remplacé, système)~; et un protocole de contrôle de dix messages pilote la plateforme (événements système, connexions et déconnexions, phases, bandes,
indicatives, croisements, heures limites, échéances, cours de référence). \texttt{PROTOCOL.md} donne toutes les règles~; la spécification TotalView-ITCH~5.0 de Nasdaq sert de modèle pour le
flux (entiers en gros-boutiste, un locate de valeur à la même place dans chaque message, horodatages en nanosecondes depuis minuit).

Une journée de client, en processus (\texttt{mx\_exchsim.scenario}), montre les rapports~: un achat de 300 à 100.00 est accepté~; une vente de 100 à 100.00 issue de la
même session s'apparie avec lui (la \omterm{def:mx:build-a-matching-engine-and-exchange-simulator:stp}{prévention de l'auto-appariement} est désactivée pour cette session) et les deux côtés reçoivent une exécution~; une vente de 500 à
99.99 en immédiat-ou-annulation exécute 200 contre le reste, et ses 300 restants sont annulés avec le motif «~I~»~; un achat de 200 à 100.02 en post-only se repose dans le carnet~; une vente de 100 à
100.01 s'exécute contre lui à 100.02~; un achat à 100.015 est rejeté car hors de la grille de cotation («~X~»)~; une annulation de la vente déjà exécutée est rejetée car trop
tardive («~L~»)~; en fin de journée, les 100 dernières actions de l'ordre post-only expirent («~E~»).

\section{Le moteur d'appariement}

Le moteur (\texttt{firm\_exchsim\_engine}) tient un carnet par instrument (le \texttt{firm.lob} du chapitre~1) et applique les types d'ordres du chapitre~2 par
\texttt{firm.ordertypes}~: \omterm{def:mx:the-limit-order-book:orders}{ordres à cours limité} et au marché, immédiat-ou-annulation et tout-ou-rien, post-only, \omterm{def:mx:order-types-and-their-uses:hidden}{ordres cachés} et \omterm{def:mx:order-types-and-their-uses:hidden}{ordres iceberg}, pegs au prix milieu et au prix primaire,
\omterm{def:mx:order-types-and-their-uses:stop}{ordres stop}, quantités minimales, et \omterm{def:mx:build-a-matching-engine-and-exchange-simulator:stp}{prévention de l'auto-appariement}.

\begin{definition}[Prévention de l'auto-appariement, annulation à la déconnexion]\label{def:mx:build-a-matching-engine-and-exchange-simulator:stp}
La \emph{prévention de l'auto-appariement}\index{prévention de l'auto-appariement} (self-trade prevention) empêche deux ordres d'un même participant (ou groupe) de s'apparier entre eux, en annulant
l'ordre au repos, l'ordre entrant, les deux, ou en décrémentant leurs quantités. L'\emph{annulation à la déconnexion}\index{annulation à la déconnexion} (cancel on disconnect) annule les ordres au repos d'une session
lorsque celle-ci est perdue.
\end{definition}

L'appariement donne d'abord la priorité au prix, puis répartit un \omterm{def:mx:the-limit-order-book:tick}{niveau de prix} selon sa règle~: premier entré, premier sorti (priorité prix-temps, livre~1, chapitre~19), ou au prorata avec une
part optionnelle réservée au premier ordre (\texttt{firm.match} du livre~1)~; il s'arrête à une bande de prix (chapitre~25) et à la limite de l'ordre.

\omcode{code/firm/exchsim/firm_exchsim_engine.py}{331}{355}{La répartition d'un niveau de prix dans le moteur~: les ordres au repos éligibles (quantités minimales respectées), exécutés dans l'ordre temporel pour la priorité prix-temps ou selon la règle au prorata configurable du livre~1, chaque exécution se faisant au prix de l'ordre au repos.}\label{lst:mx:build-a-matching-engine-and-exchange-simulator:match}

Les phases proviennent d'un calendrier~: pré-ouverture et call d'ouverture, négociation en continu ou \omterm{def:mx:auction-mechanics-and-theory:fba}{enchères par lots fréquentes} (chapitre~10), call et enchère de clôture, suspensions et
pauses avec appels de réouverture (les politiques de \texttt{firm.halts} du chapitre~25 agissent par les mêmes contrôles), le tout croisé par l'enchère du chapitre~10
(\texttt{firm.auctionsim}). Des limiteurs bornent le débit de messages de chaque session par un seau à jetons~; les frais sont prélevés par exécution selon l'indicateur de liquidité.

\section{Le flux de données de marché~: séquencement, redondance, reprise}

\begin{definition}[Lignes de flux redondantes, reprise par instantané, demande de retransmission]\label{def:mx:build-a-matching-engine-and-exchange-simulator:feed}
Les \emph{lignes de flux redondantes}\index{lignes de flux redondantes} (redundant feed lines) transportent deux fois les mêmes paquets séquencés, par des chemins distincts, afin qu'un récepteur puisse prendre chaque paquet
sur la ligne qui le livre en premier. La \emph{reprise par instantané}\index{reprise par instantané} (snapshot recovery) reconstruit un carnet à partir d'un instantané périodique de son état, étiqueté
par le dernier numéro de séquence incrémental qu'il inclut, suivi des messages incrémentaux qui le suivent. Une \emph{demande de retransmission}\index{demande de retransmission} (retransmission request) s'adresse au
service de reprise de la plateforme pour obtenir des messages par numéro de séquence.
\end{definition}

Le flux circule en paquets MoldUDP64 (un nom de session, le numéro de séquence du premier message, un nombre de messages~; un nombre nul est un battement de cœur) sur deux
lignes, A et B, chacune avec ses propres perturbations tirées de flux aléatoires issus de générateurs à compteur~: pertes indépendantes, rafales de pertes par un canal Gilbert--Elliott
à deux états, duplication, gigue, pannes. La documentation MDP~3.0 du CME décrit la même pratique~: traiter les deux flux, prendre chaque numéro de séquence
sur celui qui arrive en premier, et traiter un trou sur les deux comme une perte qui exige une reprise~; ses messages d'instantané portent le dernier numéro de séquence incrémental traité
afin qu'un récepteur puisse raccorder les deux.

Dix minutes du flux synthétique du livre~7 passées par la plateforme (\texttt{mx\_exchsim.recovery}), avec des rafales de pertes tirées indépendamment sur chaque ligne et une
panne de dix secondes sur la ligne B~: sur 5\,793 messages, la ligne A en a perdu 308 et la ligne B 517 (166 d'entre eux pendant la panne)~; en prenant chaque message sur l'une ou l'autre ligne, il en manquait encore 30, et 30 \omterm{def:mx:build-a-matching-engine-and-exchange-simulator:feed}{demandes de retransmission} les ont comblés,
de sorte que le récepteur a eu chaque message. Un instantané pris au milieu de la journée (la plateforme en prenait un toutes les 30 secondes), plus les messages incrémentaux qui le suivent,
a reconstruit un carnet identique au moteur, dix niveaux de profondeur de chaque côté.

\omcode{code/microstructure/26-build-a-matching-engine-and-exchange-simulator/python/mx_exchsim.py}{118}{123}{Les restes de l'arbitrage entre lignes~: les numéros de séquence qu'aucune ligne n'a livrés, demandés chacun au service de retransmission, dont le paquet MoldUDP64 porte le message manquant.}\label{lst:mx:build-a-matching-engine-and-exchange-simulator:gaps}

\section{Le protocole de saisie des ordres et ses sessions}

Les participants se connectent par des sessions~: un agent en processus en a une par plateforme, avec son propre modèle de latence (délais de saisie, d'accusé de réception et de données, gigue,
pics)~; un client en direct se connecte au serveur C++20 de cette machine par SoupBinTCP, le protocole de session point à point de Nasdaq au-dessus de TCP (connexion avec un
numéro de séquence demandé, données séquencées, battements de cœur), et reçoit le flux par multidiffusion en boucle locale. Les sessions portent les règles de la plateforme~: \omterm{def:mx:build-a-matching-engine-and-exchange-simulator:stp}{annulation
à la déconnexion}, groupes d'auto-appariement, limiteurs, et identifiants d'ordre client numérotés par session (un agent multi-plateforme doit rendre les siens uniques entre plateformes,
comme le routeur du chapitre~18). Les \omterm{def:mx:build-a-matching-engine-and-exchange-simulator:oe}{rapports d'exécution} ne vont qu'à la session qui possède l'ordre~; le flux va à tout le monde, et une copie de drop à qui la demande.

\section{Horloge, latence et déterminisme}

Le simulateur est une boucle à événements discrets sur une horloge simulée (livre~7, chapitre~17)~: chaque message, minuterie et contrôle est un événement à un instant en nanosecondes, et les égalités sont départagées
par un ordre fixe des classes d'événements et par l'ordre d'insertion. Chaque nombre aléatoire (latences, perturbations du flux, flux de fond) provient d'un
générateur à compteur (livre~4, chapitre~26) indexé par la graine et par son usage, de sorte qu'ajouter un participant ne modifie pas
les nombres aléatoires que voient les autres. Relancer deux fois la même journée de deux minutes donne les mêmes octets~: le SHA-256 de l'enregistrement de la ligne A et celui du
journal du moteur coïncident (\texttt{f7ede7cfa7b3\dots}).

Le moteur existe en trois exemplaires~: la référence Python, un moteur C++20 et un moteur Rust. Un jeu d'essai commun (quatre minutes de flux de bande, plus un agent chaotique
qui envoie chaque type de message et chaque contrôle~: 17\,059 enregistrements de journal) est rejoué par chacun, et chacun doit reproduire la sortie du moteur Python, les messages de flux,
rapports et instantanés finaux, octet par octet. Mesuré sur un ordinateur portable (Intel Core Ultra 7 155H) sous WSL2, sans cœurs isolés, le moteur Python
traite le jeu d'essai à environ 60\,000 enregistrements par seconde~; le programme de test C++20 (qui décode aussi et vérifie 4\,866 messages de référence, et compare chaque octet)
le parcourt à plus de 500\,000 par seconde. Le serveur C++20 en direct sert le même moteur aux clients en local.

\section{Tutoriel~: journée de certification}\label{tut:mx:build-a-matching-engine-and-exchange-simulator:tut}

\begin{tutorial}
\textbf{Objectif.} Négocier sur le simulateur en processus et par son serveur, reprendre un flux avec pertes, et certifier les trois moteurs.
\textbf{État final~:} les nombres des sections 1, 3 et 5.
\begin{tutsteps}
\item \textbf{En processus.} \texttt{Simulator(ExchangeConfig)}, un \texttt{Agent} avec \texttt{on\_report}~; \texttt{mx\_exchsim.scenario()}.
\item \textbf{Flux.} \texttt{FeedConfig(burst\_loss, outages, snapshot\_every\_ns)}~; \texttt{res.recorded(line)}, \texttt{res.retransmit(venue, seq, count)},
les instantanés~; \texttt{recovery()}.
\item \textbf{Déterminisme et identité.} \texttt{determinism()}~; \texttt{code/firm/exchsim/cpp/exchsim\_engine\_test.cpp} et \texttt{cargo test} dans
\texttt{code/firm/exchsim/rust} sur le jeu d'essai~; \texttt{throughput()}.
\item \textbf{En direct.} \texttt{code/firm/exchsim/cpp/build.sh}, puis \texttt{bin/exchsim\_server --config cfg.json} et un client SoupBinTCP
(\texttt{code/firm/exchsim/tests/test\_serve.py} en montre un en Python).
\end{tutsteps}
\textbf{Pour aller plus loin.} Ajouter l'appariement implicite des calendar spreads (le chapitre~21 le tarifait hors du moteur)~; alimenter le simulateur avec les agents du chapitre~27
en arrière-plan~; certifier un quatrième moteur par rapport au jeu d'essai.
\end{tutorial}

\section{Construction~: simulateur de bourse}\label{bld:mx:build-a-matching-engine-and-exchange-simulator:exchsim}

\begin{build}
\bfield{Rôle} Le lieu de négociation sous-jacent à chaque marché simulé des livres~10 à 13~: appariement, phases et enchères, suspensions, frais, flux, sessions.
\bfield{Interface}
\texttt{ExchangeConfig(venue, instruments, fees, phases, halts, throttle, engine\_ns, speed\_bump\_ns, batch\_interval\_ns, feed, faults, reference)}, \texttt{Simulator(venues, sip, seed)} avec \texttt{add\_background}, \texttt{add\_agent}, \texttt{add\_events},
\texttt{schedule\_call}, \texttt{run}~; \texttt{Result} avec \texttt{tape}, \texttt{feed\_messages}, \texttt{recorded}, \texttt{retransmit},
\texttt{snapshot\_bytes}, \texttt{reports}, \texttt{sip}~; préréglages~; le moteur, le codec et le serveur en C++20~; le moteur et le codec en Rust (contrat~:
\texttt{code/firm/INTERFACES.md}, section~7).
\bfield{Règles} Messages binaires en gros-boutiste générés à partir d'une seule table~; prix exprimés en 1/10\,000 d'unité monétaire~; aléa à compteur indexé par finalité~;
aucune horloge murale nulle part dans le moteur.
\bfield{Tests d'acceptation} \texttt{code/firm/exchsim/tests/}~: treize scénarios de règles du moteur~; allers-retours du codec et du schéma~; décodeur du livre~1 sur le
flux~; déterminisme par empreinte~; le moteur Python face au SHA-256 du jeu de référence~; séquencement, pertes et reprise~; reconstruction à partir d'un instantané~; agents et export de la bande~;
flux consolidé et plateformes obscures~; les moteurs C++20 et Rust identiques octet par octet sur le jeu de référence~; allers-retours des serveurs Python et C++20.
\bfield{Pour aller plus loin} Spreads implicites~; une passerelle FIX~; une seconde disposition de flux en Simple Binary Encoding.
\end{build}

\begin{omsources}
\item Nasdaq, spécification TotalView-ITCH 5.0 ; spécification du protocole MoldUDP64 ; spécification de SoupBinTCP.
\item CME Group, MDP~3.0: incremental feed arbitration; market data snapshot full recovery (client systems wiki).
\item FIX Trading Community, Simple Binary Encoding.
\item E.~N.~Gilbert, ``Capacity of a burst-noise channel'', \emph{Bell System Technical Journal} 39(5), 1960.
\end{omsources}

\section{Exercices}

\begin{exercise}[$\star$]\label{exo:mx:build-a-matching-engine-and-exchange-simulator:1}
Un paquet MoldUDP64 porte le numéro de séquence 5\,001 et un compte de 12. Quel numéro de séquence porte le paquet de données suivant~?
\end{exercise}

\begin{exercise}[$\star$]\label{exo:mx:build-a-matching-engine-and-exchange-simulator:2}
Dans le run de reprise, quelle part des messages chaque ligne a-t-elle perdue, et quelle part les deux lignes ont-elles perdue en commun~?
\end{exercise}

\begin{exercise}[$\star$]\label{exo:mx:build-a-matching-engine-and-exchange-simulator:3}
Dans le scénario, pourquoi l'annulation à 7 secondes a-t-elle été rejetée, et quel motif recevrait l'annulation d'un ordre inconnu~?
\end{exercise}

\begin{exercise}[$\star\star$]\label{exo:mx:build-a-matching-engine-and-exchange-simulator:4}
Pourquoi chaque flux aléatoire doit-il être indexé par sa finalité, et non tiré d'un unique générateur dans l'ordre des événements~?
\end{exercise}

\begin{exercise}[$\star\star$]\label{exo:mx:build-a-matching-engine-and-exchange-simulator:5}
Pourquoi la \omterm{def:mx:build-a-matching-engine-and-exchange-simulator:feed}{reprise par instantané} a-t-elle besoin du dernier numéro de séquence de l'instantané, et que se passe-t-il sans lui~?
\end{exercise}

\begin{exercise}[$\star\star$]\label{exo:mx:build-a-matching-engine-and-exchange-simulator:6}
Que prouve l'identité octet par octet de trois moteurs sur un jeu de référence, et que ne prouve-t-elle pas~?
\end{exercise}

\begin{exercise}[$\star\star\star$]\label{exo:mx:build-a-matching-engine-and-exchange-simulator:7}
\emph{Programmation.} Lancer la reprise avec l'interruption de la ligne B mais sans perte en rafale sur aucune des deux lignes. Combien de messages chaque ligne perd-elle, et combien de
\omterm{def:mx:build-a-matching-engine-and-exchange-simulator:feed}{demandes de retransmission} sont-elles nécessaires~?
\end{exercise}

\begin{exercise}[$\star\star\star$]\label{exo:mx:build-a-matching-engine-and-exchange-simulator:8}
\emph{Trouver la faille.} «~Le moteur C++ est huit fois plus rapide que le moteur Python, donc les sessions du simulateur tourneront huit fois plus vite en C++.~»
\end{exercise}

\section{Problème~: Jour de certification}

\begin{problem}[{Problème du week-end --- jour de certification}]\label{pb:mx:build-a-matching-engine-and-exchange-simulator:1}
Avant que d'autres équipes s'appuient sur le simulateur, il doit être certifié~: règles, reprise, déterminisme, identité entre implémentations et vitesse.

\medskip\noindent\textbf{Partie I --- Messages.}
\begin{enumerate}
\item Définir un \omterm{def:mx:build-a-matching-engine-and-exchange-simulator:sim}{simulateur de bourse}, un \omterm{def:mx:build-a-matching-engine-and-exchange-simulator:oe}{protocole de saisie des ordres}, une session et un \omterm{def:mx:build-a-matching-engine-and-exchange-simulator:oe}{rapport d'exécution}.
\item Énumérer les quatre protocoles et leurs nombres de messages.
\item Parcourir les rapports du scénario.
\item Définir la \omterm{def:mx:build-a-matching-engine-and-exchange-simulator:stp}{prévention de l'auto-appariement} et l'\omterm{def:mx:build-a-matching-engine-and-exchange-simulator:stp}{annulation à la déconnexion}.
\end{enumerate}

\medskip\noindent\textbf{Partie II --- Le moteur.}
\begin{enumerate}[resume]
\item Comment le moteur alloue-t-il un \omterm{def:mx:the-limit-order-book:tick}{niveau de prix} selon chaque règle~?
\item Qu'est-ce qui arrête l'appariement~?
\item Dans quelles phases la plateforme peut-elle se trouver, et qu'est-ce qui désapparie une enchère~?
\item Comment les politiques \texttt{firm.halts} agissent-elles sur la plateforme~?
\end{enumerate}

\medskip\noindent\textbf{Partie III --- Le flux.}
\begin{enumerate}[resume]
\item Décrire un paquet MoldUDP64 et les deux lignes.
\item Comment les pertes et les interruptions sont-elles tirées~?
\item Donner les pertes du run de reprise, les trous après arbitrage et les retransmissions.
\item Comment le carnet a-t-il été reconstruit à partir d'un instantané, et comment l'a-t-on vérifié~?
\item Que recommande la documentation MDP~3.0 de CME pour ses flux A et B~?
\end{enumerate}

\medskip\noindent\textbf{Partie IV --- Certification.}
\begin{enumerate}[resume]
\item Qu'est-ce qui rend le simulateur déterministe~?
\item Décrire le jeu de référence.
\item \emph{Énoncer le résultat nommé}~: la suite de conformité réussie par chaque moteur, l'identité octet par octet des flux des trois moteurs sur le jeu de référence, et le
nombre de messages par seconde des moteurs sur ce portable.
\item Pourquoi le chiffre du C++ est-il une borne inférieure de la vitesse du moteur~?
\item Quelles parties du simulateur ne sont pas certifiées par le jeu de référence~?
\item Que ajouteriez-vous avant qu'une équipe de trading en production l'utilise~?
\item En une phrase~: pourquoi un simulateur doit-il être déterministe~?
\end{enumerate}
\end{problem}

\section{Questions d'entretien}

\begin{interviewq}[$\star$ \iqroles{développeur}]\label{iq:mx:build-a-matching-engine-and-exchange-simulator:1}
Comment un gestionnaire de flux se rétablit-il après un trou dans un flux UDP en multidiffusion~?
\end{interviewq}

\begin{interviewq}[$\star\star$ \iqroles{développeur}]\label{iq:mx:build-a-matching-engine-and-exchange-simulator:2}
Concevriez-vous les structures de données d'un \omterm{def:mx:the-limit-order-book:engine}{moteur d'appariement} en priorité prix-temps avec des annulations rapides~?
\end{interviewq}

\begin{interviewq}[$\star\star$ \iqroles{développeur}]\label{iq:mx:build-a-matching-engine-and-exchange-simulator:3}
Comment rendriez-vous un simulateur de marché reproductible à l'octet près, aléa compris~?
\end{interviewq}

\begin{interviewq}[$\star\star$ \iqroles{trader}]\label{iq:mx:build-a-matching-engine-and-exchange-simulator:4}
Votre ordre IOC a été partiellement exécuté et le reste annulé. Quels rapports recevez-vous, et dans quel ordre~?
\end{interviewq}

\begin{interviewq}[$\star\star$ \iqroles{risque}]\label{iq:mx:build-a-matching-engine-and-exchange-simulator:5}
Pourquoi les plateformes proposent-elles l'\omterm{def:mx:build-a-matching-engine-and-exchange-simulator:stp}{annulation à la déconnexion} et la \omterm{def:mx:build-a-matching-engine-and-exchange-simulator:stp}{prévention de l'auto-appariement}, et quels risques chacune de ces protections contrôle-t-elle~?
\end{interviewq}

\begin{interviewq}[$\star\star\star$ \iqroles{développeur}]\label{iq:mx:build-a-matching-engine-and-exchange-simulator:6}
Deux implémentations du même moteur divergent sur un message sur un million. Comment découvririez-vous laquelle a tort~?
\end{interviewq}
